\section{Standing setup}
\label{a:sec:setup}

This section fixes what is carried over from \ChM.

Let $Z_n$ be the binary Galton--Watson process started from a single progenitor,
$Z_0=1$, in which each individual independently divides into two with probability
$p$ and dies with probability $1-p$ at each generational increment. Its offspring
generating function is $\phi(z) = pz^2+(1-p)$, the mean population is
$\ex{Z_n} = (2p)^n$, and the survival probability $S_n = \pr{Z_n>0}$ satisfies
\begin{equation}
\label{a:eq:SnIter}
S_{n+1} = 2pS_n - pS_n^2,\qquad S_0 = 1 .
\end{equation}
Throughout this chapter $p$ is confined to the subcritical range
\begin{equation}
\label{a:eq:range}
0\le p<\tfrac12,
\end{equation}
where extinction is certain and $S_n\downarrow0$. Two derived symbols are used
constantly:
\begin{equation}
\label{a:eq:epsr}
\varepsilon = 1-2p\in(0,1] \quad\text{(the distance from criticality)},
\qquad
r = 2p\in[0,1) \quad\text{(the logistic multiplier)} .
\end{equation}

Because the quadratic term in \cref{a:eq:SnIter} becomes negligible once $S_n$ is
small, the decay is asymptotically geometric with ratio $2p$, and the constant
$A(p)$ of \cref{a:eq:Apdef} is the amplitude of that decay:
\begin{equation}
\label{a:eq:lateS}
S_n\sim A(p)\,(2p)^n\qquad(n\to\infty) .
\end{equation}
The relation is an asymptotic equivalence, not an identity: $S_n$ itself tends to
zero, and what converges is the ratio $S_n/(2p)^n$. That this ratio converges to
something finite and strictly positive is not assumed here; it is proved in
\cref{a:sec:product}.

Splitting
$\ex{Z_n}$ over the events $\{Z_n=0\}$ and $\{Z_n>0\}$, the first contributing
nothing, gives $\ex{Z_n\mid Z_n>0} = (2p)^n/S_n$ and hence
\begin{equation}
\label{a:eq:qsmean}
\ex{Z_n\mid Z_n>0}\;\longrightarrow\;\frac{1}{A(p)},
\end{equation}
the \emph{quasi-stationary mean}; and the corresponding second-moment calculation
gives a limiting conditional variance
\begin{equation}
\label{a:eq:condvar}
\var{Z_n\mid Z_n>0}\;\longrightarrow\;
\frac{1}{A(p)}\lrb{1+\frac{1}{1-2p}} - \frac{1}{A(p)^2} ,
\end{equation}
again a function of $p$ and $A(p)$ alone. The derivations of
\cref{a:eq:qsmean,a:eq:condvar}, together with the existence of the Yaglom limit
they describe, are given in \ChM{} and are not repeated.

One comparison object is also inherited. The continuous-time linear birth--death
process with per-capita birth rate $\lambda$ and death rate $\mu>\lambda$ has a
limiting conditional mean $\mu/(\mu-\lambda)$, and matching the parametrisations
through the competing-clocks identity $p=\lambda/(\lambda+\mu)$ --- proved in
\ChM{} and used here by citation --- turns this into
\begin{equation}
\label{a:eq:Acts}
\Ac(p) = \frac{1-2p}{1-p} = \frac{2\varepsilon}{1+\varepsilon},
\qquad
\frac{1}{\Ac(p)} = \frac{1-p}{1-2p} .
\end{equation}
Unlike $A(p)$, the continuous-time constant $\Ac(p)$ is elementary.

Finally, an endpoint convention. At $p=0$ the process dies at the first step and
\cref{a:eq:SnIter} degenerates; we set
\begin{equation}
\label{a:eq:endpoint}
A(0) := \lim_{p\downarrow0}A(p) = \tfrac12,
\end{equation}
a value confirmed independently by the product of \cref{a:sec:product} and by the
parity bound of \cref{a:prop:parity}.
